\documentclass[twocolumn]{aastex631}
\usepackage{amsmath}
\usepackage{booktabs}
\shorttitle{SDSS denominator/proxy atlas}
\shortauthors{NebulaMind}
\begin{document}

\title{Supplementary SDSS Denominator and Proxy Atlas for Galaxy-Evolution Follow-up}
\author{NebulaMind Research Autopilot}
\affiliation{Public SDSS DR17 data only}

\begin{abstract}
This supplement collects eight SDSS DR17 denominator and proxy notes that share the same capped 60k-row pilot cache and the same selection-function caveats. The 60,000-row cache is an arbitrary, non-random pilot-query cap, not a physical or volume-limited selection effect, so all counts and fractions remain conditional on the SDSS optical selection used here. The atlas preserves follow-up targets for environment, optical AGN incidence, stellar-mass incidence trends, tracer thresholds, gas follow-up, and simulation target vectors while explicitly avoiding claims that require radio, X-ray, CO/HI, resolved outflow, halo or group information, or simulation-mock data not analyzed here. These counts and fractions are conditional on the SDSS optical selection used here, not global volume-limited statistics, and the cached coverage is 24.0\% of the strict four-line S/N$\geq3$ parent. It is one follow-up atlas, not eight independent causal-feedback papers. Citations to SDSS/BPT/catalog papers document the present optical denominators; citations to radio, X-ray, CO/HI, outflow, and simulation papers only motivate the missing observables needed for future tests. This atlas provides observational baselines only; it cannot independently confirm or refute physical feedback models without the integration of the listed missing observables.
\end{abstract}

\keywords{galaxies: evolution --- surveys --- catalogs --- methods: observational --- methods: statistical}

\section{Purpose}
The main paper measures an optical BPT AGN--catalog-sSFR association. These eight topics are distinct: each is scientifically useful as a denominator or target-vector definition, but each lacks at least one core physical observable required by its original proposal. The BPT language and catalog-backbone language here follow the same SDSS/MPA-JHU-style value-added tables and standard demarcations as the flagship \citep{sdssdr17,brinchmann2004,york2000,baldwin1981,kewley2001,kauffmann2003bpt,kewley2006,stasinska2008,stasinska2015}. The radio/X-ray/CO/HI/outflow/simulation references that appear later in the notes are therefore role-separated as future-data motivation, not validation of the present optical denominators. Keeping them in one supplement prevents overclaiming and gives future work a single checklist of what still must be added.

\section{Shared denominator}
The atlas uses the same cached public-data backbone as the main paper: 60,000 cached rows from a strict public four-line S/N$\geq3$ parent of 249,917 rows, i.e. 24.0\% cached coverage. The four-line selection is sSFR-dependent and the cache is capped and non-random, so all counts and fractions are conditional denominators rather than population-complete measurements. The row-level stellar masses and catalog sSFR values are taken from the public MPA-JHU-style \texttt{galSpecExtra} table after the same SDSS joins used in the flagship \citep{sdssdr17,brinchmann2004,york2000}. The SDSS/BPT/catalog references support these observed denominators; the later multiwavelength and simulation references only mark the follow-up measurements that are still missing.

The eight subsections below are intentionally parallel: each one states the observed optical denominator or target vector, then lists the missing observables that a future multiwavelength or simulation-based test would have to add before any physical inference can be made.

\begin{deluxetable*}{lrrr}
\tabletypesize{\scriptsize}
\tablecaption{Selection cascade shared by the atlas.\label{tab:supp-selection}}
\tablehead{\colhead{Selection stage} & \colhead{Public DR17 rows} & \colhead{Cached rows} & \colhead{Retention vs. spectro-z parent}}
\startdata
SpecObj GALAXY, 0.02<z<0.12 & 501,060 & -- & 100.0\% \\
plus galSpecInfo/PhotoObj/galSpecExtra and mass/sSFR bounds & 416,554 & -- & 83.1\% \\
plus galSpecLine join & 416,554 & -- & 83.1\% \\
four BPT lines positive with positive errors & 373,445 & 60,000 & 74.5\% \\
four BPT lines S/N>=3 & 249,917 & 60,000 & 49.9\% \\
four BPT lines S/N>=5 & 176,523 & 42,446 & 35.2\% \\
four BPT lines S/N>=10 & 91,768 & 22,311 & 18.3\% \\
\enddata
\tablecomments{Counts are read-only public SDSS DR17 count queries plus the cached analysis table. Retention is shown as a percentage of the spectro-z parent. Cached rows are shown only where the cache applies. The sharp retention drop at higher S/N mainly reflects preferential loss of passive galaxies from the emission-line denominator, so the surviving cache becomes less representative of quiescent hosts as the cut tightens.}
\end{deluxetable*}

\section{Atlas notes}

\subsection{Environment baseline: SDSS density proxy for low-sSFR incidence}
This note isolates an internal environmental denominator that can later be joined to group catalogs and halo masses. Within this selection-biased emission-line denominator, the relative 10th-neighbor index covaries with the catalog low-sSFR fraction; this index is only an internal ordinal rank and does not map to physical environmental volume density or halo density. The SDSS emission-line denominator contains 60,000 galaxies with an internally computed 10th-neighbor index. The high-density quartile has a low-sSFR emission-line fraction of 0.230 (3,456/15,000), while the low-density quartile has 0.181 (2,710/15,000). The bootstrap high-minus-low interval is [0.041, 0.059], and a linear probability model adjusted for log stellar mass and redshift gives a high-density coefficient of 0.032 +/- 0.004. This is a denominator-level environmental diagnostic; the missing observables are:
\begin{itemize}
\item group catalogues
\item robust central/satellite labels
\item halo masses
\item spectroscopic fiber-collision correction at the 55-arcsec scale
\item morphology
\item multi-redshift selection functions
\end{itemize}
Within this selection-biased emission-line cache, the 10th-neighbor statistic is only a relative local rank, not a physical volume density and not a substitute for central/satellite labels or a volume-complete halo-density measurement.
These are still needed for a future environmental test \citep{peng2010,wetzel2013,dekel2006}.

\begin{figure}
\centering
\includegraphics[width=\columnwidth]{../figures/topic-01.pdf}
\caption{SDSS optical emission-line denominator: the low-sSFR emission-line fraction as a function of the 10th-neighbor density proxy in the SDSS emission-line sample. This is a selection-dependent baseline for future group- and halo-matched follow-up, not a physical-feedback measurement.}
\label{fig:m1-rp2-environment-quenching}
\end{figure}


\subsection{Maintenance-heating denominator: optical AGN in massive SDSS hosts}
This note identifies the optical-AGN duty-cycle denominator that radio and X-ray data would need to test maintenance heating. Among massive, low-sSFR SDSS emission-line galaxies, the optical AGN fraction can be used as a denominator for X-ray and radio maintenance-heating follow-up. The massive subset (\(\log M_\star \geq 10.8\)) contains 9,298 emission-line galaxies, of which 5,695 are low-sSFR by the pilot threshold. The optical BPT AGN fraction is 0.430 in the massive subset and 0.607 among massive low-sSFR objects. This provides an optical duty-cycle denominator for X-ray and radio follow-up, not a heating-to-cooling measurement. The missing observables are:
\begin{itemize}
\item X-ray cavity or cooling-luminosity measurements
\item radio jet powers
\item halo-selected parent catalogues
\item nondetection modelling
\end{itemize}
These are still needed for a future maintenance-heating test \citep{best2005,heckmanbest2014,fabian2012,mcnamara2007,lamassa2013}.

\begin{figure}
\centering
\includegraphics[width=\columnwidth]{../figures/topic-02.pdf}
\caption{SDSS optical emission-line denominator: the massive and low-sSFR SDSS emission-line subsets used as a baseline for future X-ray and radio measurements, not a heating-to-cooling result.}
\label{fig:m1-rp3-maintenance-heating}
\end{figure}


\subsection{Outflow-kinematics denominator: high-excitation SDSS AGN}
This note isolates the high-excitation optical-AGN denominator that resolved kinematics would need to test escape versus recycling. High-excitation optical AGN candidates number 4,440 of 60,000 emission-line galaxies (0.074). Their median \(\log {\rm sSFR}\) is -11.53, compared with -10.14 for the full denominator. SDSS does not measure escape velocity or multiphase outflow velocities here; the note supplies a denominator for resolved follow-up rather than an escape or recycling result. The missing observables are:
\begin{itemize}
\item resolved outflow velocities
\item halo potentials
\item molecular, ionized, and neutral gas phases
\item CGM recycling tracers
\end{itemize}
These are still needed for a future outflow test \citep{veilleux2005,cicone2014,carniani2017,fiore2017,lamassa2013}.

\begin{figure}
\centering
\includegraphics[width=\columnwidth]{../figures/topic-03.pdf}
\caption{SDSS optical emission-line denominator: the high-excitation AGN subset used to define an observational baseline for future resolved-kinematic measurements, not an escape or recycling result.}
\label{fig:m2-p1-outflow-escape-recycling}
\end{figure}


\subsection{Radio-jet environment baseline: optical AGN fraction vs. density proxy in massive hosts}
This note defines the environment-stratified optical denominator that future radio and X-ray work could test. The local-density proxy is correlated with the optical AGN fraction in massive SDSS hosts and motivates environment-stratified radio and X-ray follow-up. Among massive hosts, the high-density quartile has an optical AGN fraction of 0.509, while the low-density quartile has 0.367. The bootstrap high-minus-low interval is [0.112, 0.170]. This is an optical/environment denominator for future radio-jet follow-up; it does not measure radio jet power or coupling efficiency. The missing observables are:
\begin{itemize}
\item radio jet morphology and age
\item cavity or shock energetics
\item hot-gas density
\item calibrated jet-power estimates
\end{itemize}
These are still needed for a future radio-jet test \citep{best2005,mcnamara2007,heckmanbest2014}.

\begin{figure}
\centering
\includegraphics[width=\columnwidth]{../figures/topic-04.pdf}
\caption{SDSS optical emission-line denominator: the high- and low-density quartile comparison among massive SDSS hosts, used as a baseline for future radio-jet and X-ray work, not a coupling measurement.}
\label{fig:m2-p2-radio-jet-environment}
\end{figure}


\subsection{Mass-bin diagnostic: low-sSFR and optical AGN incidence}
This note pins down the mass bin where a future gas-inclusive study should look for an incidence change. We measure the incidence of low catalog-sSFR and optical AGN classification across stellar-mass bins in this emission-line subset. What stellar-mass bin contains the highest representation of low-sSFR and optical AGN classifications within this selection-biased SDSS denominator? The first stellar-mass bin with low-sSFR fraction above 0.5 is \(\log(M_\star/M_\odot) \in [11.0,12.5]\). The optical AGN fraction peaks in the 11.0-12.5 bin at 0.520. This is an optical distribution diagnostic; gas fractions and baryon deficits are needed before assigning any physical meaning to the apparent incidence change. The missing observables are:
\begin{itemize}
\item gas fractions
\item baryon deficits
\item halo masses
\item stellar-feedback observables
\item high-redshift extensions
\end{itemize}
The same binning is therefore best treated as a population-distribution diagnostic, not a statement about a transition mass for individual galaxies \citep{peng2010,wetzel2013,dekel2006}. In this optical-emission-line denominator, the 11.0--12.5 dex peak is a selection-function artifact: the S/N$\geq$3 cut preferentially removes truly passive, massive galaxies, leaving a surviving emission-line subset that is artificially concentrated in that mass bin. It must not be interpreted as a universal feedback threshold.

\begin{figure}
\centering
\includegraphics[width=\columnwidth]{../figures/topic-05.pdf}
\caption{SDSS optical emission-line denominator: mass-bin diagnostic for low-sSFR and optical AGN incidence in the SDSS emission-line denominator. This is a population baseline for future gas-inclusive follow-up, not a physical transition-mass measurement. The 11.0--12.5 dex peak is a selection-function artifact in this emission-line cache, not a universal feedback threshold.}
\label{fig:m2-p3-feedback-transition-mass}
\end{figure}


\subsection{Tracer-threshold census for multiphase follow-up}
This note compares optical tracer choices against one shared denominator before any multiphase census is attempted. How strongly do simple optical tracer definitions change the inferred AGN or feedback-candidate prevalence in one common SDSS denominator? Within the same 60,000-galaxy denominator, simple optical tracer definitions produce prevalence from 0.136 to 0.418. The widest-to-narrowest prevalence ratio is 3.1 before adding molecular, neutral, or X-ray or radio phases. This demonstrates why a common-denominator multiphase census is required; it does not measure molecular or neutral outflow rates. The missing observables are:
\begin{itemize}
\item ionized, molecular, and neutral tracers
\item X-ray or radio tracers
\item a shared parent denominator
\item a consistent aperture model
\end{itemize}
These are still needed for a future multiphase test \citep{xcoldgass2017,xgass2018,cicone2014,carniani2017,fiore2017,veilleux2005}.

\begin{figure}
\centering
\includegraphics[width=\columnwidth]{../figures/topic-06.pdf}
\caption{SDSS optical emission-line denominator: prevalence of alternative tracer definitions within the 60,000-galaxy sample. This is a baseline for future multiphase work, not a molecular or neutral gas census.}
\label{fig:m3-p1-multiphase-census}
\end{figure}


\subsection{Gas-depletion denominator: optical baseline for CO/HI follow-up}
This note defines the denominator for CO/HI gas-fraction and depletion-time follow-up. How many massive low-sSFR or transitioning SDSS galaxies with valid emission-line measurements are available as a denominator for CO gas-fraction and depletion-time follow-up? The massive low-sSFR denominator contains 6,729 galaxies in the SDSS emission-line sample. Its optical BPT AGN fraction is 0.549, and the median H-alpha luminosity proxy is 40.06. Here the H-alpha luminosity proxy is the aperture-corrected \texttt{galSpecExtra} catalog value, not raw fiber flux. The median H-alpha luminosity proxy is 0.66 dex lower than in massive star-forming emission-line galaxies. SDSS optical data cannot distinguish molecular-gas depletion from reduced star-formation efficiency; this note identifies the CO/HI follow-up denominator and optical baseline. The missing observables are:
\begin{itemize}
\item CO or dust-based molecular gas masses
\item aperture-matched SFRs
\item morphology
\item environment labels
\end{itemize}
These are still needed for a future gas-fraction or depletion-time test \citep{xcoldgass2017,xgass2018,piotrowska2022}.

\begin{figure}
\centering
\includegraphics[width=\columnwidth]{../figures/topic-07.pdf}
\caption{SDSS optical emission-line denominator: the massive low-sSFR SDSS galaxies available for CO/HI depletion-time follow-up, not a gas-depletion-efficiency measurement.}
\label{fig:m3-p2-gas-depletion-efficiency}
\end{figure}


\subsection{Simulation target vector for forward-model comparison}
This note provides a compact observed target vector for forward modelling, not a direct simulation comparison. What compact SDSS target vector of low-sSFR fraction, optical AGN incidence, and colour versus mass and redshift can be used for forward-model comparison? The pilot writes 15 mass-redshift cells with \(n \geq 50\) as a compact comparison vector. Across mass bins, low-sSFR fractions span 0.005-0.729, and optical AGN fractions span 0.003-0.520. The output is an observed target vector for simulation forward modelling, not a direct simulation comparison. The missing observables are:
\begin{itemize}
\item simulation mocks passed through the same optical S/N and fiber-aperture selection function used here, then through the SDSS, MaNGA, ALMA, X-ray, and radio selection functions
\item aperture models
\item noise models
\end{itemize}
Without those matched selection steps, any simulation comparison is not a valid test. These are still needed for a future simulation-comparison test \citep{simba2019,tng2019,eagle2015}.

\begin{figure}
\centering
\includegraphics[width=\columnwidth]{../figures/topic-08.pdf}
\caption{SDSS optical emission-line denominator: low-sSFR fraction, optical AGN incidence, and colour versus mass and redshift in the SDSS emission-line sample. This is an observed target vector for forward modelling, not a direct simulation comparison.}
\label{fig:m3-p3-simulation-validation}
\end{figure}

\section{Atlas summary}
Table~\ref{tab:atlas-summary} condenses the follow-up menu across the eight notes. All eight notes are linked by the same limitation: they remain SDSS optical denominators or target vectors until the missing multiwavelength, morphological, or mock-observation data are added, so their present role is to organize follow-up rather than to establish causal physical claims.

\begin{deluxetable*}{llll}
\tabletypesize{\scriptsize}
\tablecaption{Atlas-level follow-up menu. Each row summarizes the present optical role and the missing observables needed before any physical inference.\label{tab:atlas-summary}}
\tablehead{\colhead{Topic} & \colhead{Observed baseline} & \colhead{Missing observables} & \colhead{Role}}
\startdata
Environment & low-sSFR vs.\ 10th-neighbor rank & group catalogs; central/satellite labels; halo mass; fiber-collision correction & environment test \\
Maintenance heating & optical AGN in massive low-sSFR hosts & X-ray cavities; cooling luminosity; radio jet powers; halo-selected parents & radio/X-ray follow-up \\
Outflow kinematics & high-excitation AGN subset & resolved velocities; halo potentials; multiphase gas; CGM tracers & kinematic follow-up \\
Env.\ jets & density-stratified AGN fraction & radio morphology/age; cavity energetics; hot-gas density & radio-jet follow-up \\
Mass bin & low-sSFR and AGN by $M_\star$ bin & gas fractions; baryon deficits; halo masses; feedback observables & selection diagnostic \\
Tracer census & tracer prevalence in 60k sample & multiphase tracers; shared denominator; aperture model & multiphase follow-up \\
Gas depletion & massive low-sSFR baseline; H$\alpha$ proxy & CO/dust gas masses; aperture-matched SFRs; morphology; environment & CO/HI follow-up \\
Simulation vector & mass-redshift target vector & mocks through SDSS/MaNGA/ALMA/X-ray/radio selection; aperture/noise models & forward model \\
\enddata
\tablecomments{The table is a compact index of the subsection-level missing-observables lists; it does not add new measurements or change any counts. The sharp retention drop at higher S/N mainly reflects preferential loss of passive galaxies from the emission-line denominator, so the surviving cache becomes less representative of quiescent hosts as the cut tightens.}
\end{deluxetable*}


\section{Package decision}
These eight notes should remain supplementary until the missing observables are added. They are suitable as follow-up target definitions, denominator baselines, or appendix material under the main result, but not as independent causal-feedback papers in their current SDSS-only form.

\section{Local reproducibility}
This PDF was generated from the local candidate package \texttt{RP1\_FLAGSHIP\_WITH\_SUPPLEMENT\_20260709T013510Z}. It does not replace any public-linked PDF and does not touch public pages, live roots, product databases, deployment state, git history, billing/OAuth state, cron jobs, or external submission systems.


\begin{thebibliography}{}
\bibitem[Abdurro'uf et al.(2022)]{sdssdr17} Abdurro'uf, Accetta, K., Aerts, C., et al. 2022, ApJS, 259, 35
\bibitem[Baldwin et al.(1981)]{baldwin1981} Baldwin, J.~A., Phillips, M.~M., \& Terlevich, R. 1981, PASP, 93, 5
\bibitem[Best et al.(2005)]{best2005} Best, P.~N., Kauffmann, G., Heckman, T.~M., et al. 2005, MNRAS, 362, 25
\bibitem[Brinchmann et al.(2004)]{brinchmann2004} Brinchmann, J., Charlot, S., White, S.~D.~M., et al. 2004, MNRAS, 351, 1151
\bibitem[Carniani et al.(2017)]{carniani2017} Carniani, S., Marconi, A., Maiolino, R., et al. 2017, A\&A, 605, A42
\bibitem[Catinella et al.(2018)]{xgass2018} Catinella, B., Saintonge, A., Janowiecki, S., et al. 2018, MNRAS, 476, 875
\bibitem[Cicone et al.(2014)]{cicone2014} Cicone, C., Maiolino, R., Sturm, E., et al. 2014, A\&A, 562, A21
\bibitem[Dave et al.(2019)]{simba2019} Dave, R., Angles-Alcazar, D., Narayanan, D., et al. 2019, MNRAS, 486, 2827
\bibitem[Dekel \& Birnboim(2006)]{dekel2006} Dekel, A., \& Birnboim, Y. 2006, MNRAS, 368, 2
\bibitem[Fabian(2012)]{fabian2012} Fabian, A.~C. 2012, ARA\&A, 50, 455
\bibitem[Fiore et al.(2017)]{fiore2017} Fiore, F., Feruglio, C., Shankar, F., et al. 2017, A\&A, 601, A143
\bibitem[Heckman \& Best(2014)]{heckmanbest2014} Heckman, T.~M., \& Best, P.~N. 2014, ARA\&A, 52, 589
\bibitem[Kauffmann et al.(2003a)]{kauffmann2003bpt} Kauffmann, G., Heckman, T.~M., Tremonti, C., et al. 2003a, MNRAS, 346, 1055
\bibitem[Kewley et al.(2001)]{kewley2001} Kewley, L.~J., Dopita, M.~A., Sutherland, R.~S., Heisler, C.~A., \& Trevena, J. 2001, ApJ, 556, 121
\bibitem[Kewley et al.(2006)]{kewley2006} Kewley, L.~J., Groves, B., Kauffmann, G., \& Heckman, T. 2006, MNRAS, 372, 961
\bibitem[LaMassa et al.(2013)]{lamassa2013} LaMassa, S.~M., Heckman, T.~M., Ptak, A., \& Urry, C.~M. 2013, ApJL, 765, L33
\bibitem[McNamara \& Nulsen(2007)]{mcnamara2007} McNamara, B.~R., \& Nulsen, P.~E.~J. 2007, ARA\&A, 45, 117
\bibitem[Nelson et al.(2019)]{tng2019} Nelson, D., Springel, V., Pillepich, A., et al. 2019, Computational Astrophysics and Cosmology, 6, 2
\bibitem[Peng et al.(2010)]{peng2010} Peng, Y.-j., Lilly, S.~J., Kovac, K., et al. 2010, ApJ, 721, 193
\bibitem[Piotrowska et al.(2022)]{piotrowska2022} Piotrowska, J.~M., Bluck, A.~F.~L., Maiolino, R., \& Peng, Y.-j. 2022, MNRAS, 512, 1052
\bibitem[Saintonge et al.(2017)]{xcoldgass2017} Saintonge, A., Catinella, B., Tacconi, L.~J., et al. 2017, ApJS, 233, 22
\bibitem[Schaye et al.(2015)]{eagle2015} Schaye, J., Crain, R.~A., Bower, R.~G., et al. 2015, MNRAS, 446, 521
\bibitem[Stasinska et al.(2008)]{stasinska2008} Stasinska, G., Asari, N.~V., Cid Fernandes, R., et al. 2008, MNRAS, 391, L29
\bibitem[Stasinska et al.(2015)]{stasinska2015} Stasinska, G., Costa Duarte, M.~V., Vale Asari, N., Cid Fernandes, R., \& Sodre, L. 2015, MNRAS, 449, 559
\bibitem[Veilleux et al.(2005)]{veilleux2005} Veilleux, S., Cecil, G., \& Bland-Hawthorn, J. 2005, ARA\&A, 43, 769
\bibitem[Wetzel et al.(2013)]{wetzel2013} Wetzel, A.~R., Tinker, J.~L., Conroy, C., \& van den Bosch, F.~C. 2013, MNRAS, 432, 336
\bibitem[York et al.(2000)]{york2000} York, D.~G., Adelman, J., Anderson, J.~E., Jr., et al. 2000, AJ, 120, 1579
\end{thebibliography}

\end{document}
